\chapter{Immunisations for Travel}

\section*{Definition and Overview}
Travel immunisation protects travellers against diseases that are present in the countries they plan to visit but uncommon or absent in Australia. The vaccines recommended depend on the destination, the type and length of stay, the season, and the traveller's age and health. Some travel vaccines, such as yellow fever, are required for entry into certain countries, while others are strongly recommended to prevent illness during and after travel. Planning is essential: some vaccines need several weeks to produce full immunity, and courses may require multiple doses. A travel medicine consultation with a GP or specialist travel clinic is the best way to ensure you are protected.

\section*{Epidemiology}
Australians undertake millions of overseas trips each year, and many travel to regions where infectious diseases such as hepatitis A, typhoid, and Japanese encephalitis are more common than in Australia. Vaccine-preventable diseases acquired while travelling are a leading cause of illness in returning travellers, and some, such as malaria, are not preventable by vaccine and require other measures. Yellow fever vaccination is a legal requirement for entry to parts of Africa and South America, and proof of vaccination is checked at borders. Rates of travel-related illness are highest among travellers who do not seek pre-travel advice, so planned immunisation is a key part of staying healthy overseas.

\section*{Aetiology and Pathophysiology}
Travel vaccines protect against diseases the traveller would not normally be exposed to at home. They work through the same immune principles as all vaccines: the immune system is exposed to a safe form of the pathogen and develops antibodies and memory cells. Because the diseases are absent or rare in Australia, natural immunity is unlikely, so vaccination is the only practical way to be protected. Some travel vaccines, such as hepatitis A, provide protection for several years or life after a course; others, such as typhoid and yellow fever, require boosters at intervals. The vaccine schedule is planned backwards from the date of departure so that immunity is at its peak when the traveller arrives.

\section*{Clinical Features}
\begin{itemize}
    \item Common travel vaccines: hepatitis A, typhoid, and yellow fever
    \item Destination-specific vaccines: Japanese encephalitis, cholera, rabies, and tick-borne encephalitis
    \item Routine vaccines may also be needed: measles, diphtheria, tetanus, whooping cough, and influenza, particularly during outbreaks
    \item Malaria is not vaccine-preventable and requires antimalarial medicine plus mosquito avoidance
    \item Some vaccines require proof of vaccination for border entry, especially yellow fever
    \item Planning should start 6--8 weeks before departure
    \item Many travel vaccines are not covered by Medicare and are purchased by the traveller
\end{itemize}

\section*{Differential Diagnosis}
The right vaccine schedule is individual, and several factors determine recommendations.
\begin{itemize}
    \item \textbf{Destination and itinerary:} urban versus rural travel, specific countries, and local outbreaks
    \item \textbf{Duration and style of travel:} longer stays, backpacking, and work in healthcare carry higher risk
    \item \textbf{Traveller's health:} age, pregnancy, immunosuppression, and chronic disease alter recommendations
    \item \textbf{Routine vaccine status:} a full routine schedule reduces the need for extra doses
    \item \textbf{Required versus recommended:} yellow fever may be required for entry; other vaccines are recommended for protection
\end{itemize}

\section*{When to Seek Medical Care}
Book a travel health consultation at least 6--8 weeks before departure, and earlier if possible, so that vaccine courses can be completed. Travellers with medical conditions, pregnant women, and parents travelling with infants should seek advice early. A person who becomes unwell during or after travel should seek medical review, mentioning their destinations to the doctor.

\textbf{Call 000 (or local emergency services) for:}
\begin{itemize}
    \item High fever with confusion, severe headache, stiff neck, or a rash during or after travel (possible serious infection)
    \item Difficulty breathing or signs of a severe allergic reaction after a vaccine
    \item Severe diarrhoea with blood, persistent vomiting, or inability to keep fluids down
    \item Yellowing of the skin or eyes with fever after travel
    \item Any illness with rapid deterioration
\end{itemize}

\section*{Diagnosis and Investigations}
\begin{itemize}
    \item \textbf{Travel risk assessment:} structured review of destination, season, duration, activities, and accommodation
    \item \textbf{Immunisation record review:} checking routine vaccine status and previous travel vaccines
    \item \textbf{Serology:} blood tests where immunity is uncertain, for example hepatitis A or measles
    \item \textbf{Pregnancy testing and medical review:} where pregnancy or chronic disease affects vaccine choice
    \item \textbf{Documentation:} the International Certificate of Vaccination or Prophylaxis for yellow fever
\end{itemize}

\section*{Management}
\begin{itemize}
    \item \textbf{Hepatitis A vaccine:} recommended for most travel to developing regions; a course gives long-lasting protection
    \item \textbf{Typhoid vaccine:} for travel to South Asia and other endemic areas, especially where hygiene is poor
    \item \textbf{Yellow fever vaccine:} required for entry to parts of Africa and South America; one dose provides long-term protection
    \item \textbf{Japanese encephalitis vaccine:} for longer stays or rural travel in Asia
    \item \textbf{Rabies pre-exposure vaccine:} for travellers at risk of animal bites where access to treatment is limited
    \item \textbf{Cholera vaccine:} for high-risk travellers to outbreak areas
    \item \textbf{Routine catch-up:} measles, influenza, and other routine vaccines brought up to date before travel
    \item \textbf{Non-vaccine protection:} antimalarials, insect repellent, food and water precautions, and travellers' diarrhoea kits
\end{itemize}

\section*{Complications}
\begin{itemize}
    \item Vaccine reactions: soreness, fever, and fatigue, which are usually mild and brief
    \item Diseases acquired without vaccination: hepatitis A, typhoid, yellow fever, and others can cause severe illness and hospitalisation
    \item Border entry problems for travellers without required yellow fever certification
    \item Costs of unplanned medical care overseas
    \item Interaction of illness with pregnancy and chronic disease
\end{itemize}

\section*{Prognosis and Outcomes}
Travel vaccination is highly effective when planned in advance, and travellers who complete recommended schedules have very low rates of the targeted diseases. Even for diseases without vaccines, such as malaria, combining precautions greatly reduces risk. The main causes of preventable travel illness are failing to seek advice, leaving planning too late, and incomplete courses. With proper planning, the vast majority of travellers return home healthy, and any illness that occurs is promptly recognised and treated.

\section*{Prevention and Self-Care}
\begin{itemize}
    \item Start travel health planning 6--8 weeks before departure
    \item Complete all recommended vaccine courses before leaving
    \item Check the Smartraveller website for destination-specific health advice
    \item Protect against mosquito bites with repellent, long clothing, and treated nets where malaria or dengue risk exists
    \item Drink safe water, eat freshly cooked food, and practise good hand hygiene
    \item Carry a basic medical kit and any regular medicines
    \item Declare your travel history to any doctor if you become unwell after returning
\end{itemize}

\section*{Sources and Further Reading}
The following reputable sources provide further information for healthcare students and patients seeking reliable, current advice about travel immunisation.
\begin{itemize}
    \item Australian Government Department of Health and Aged Care. \textit{Travel vaccination and yellow fever.} https://www.health.gov.au/
    \item Smartraveller. \textit{Health advice for travellers.} https://www.smartraveller.gov.au/
    \item Travel Health Advisory Group. \textit{Travel vaccination recommendations.} https://www.travelhealth.com.au/
    \item World Health Organization. \textit{International travel and health.} https://www.who.int/travel-advice
    \item Centers for Disease Control and Prevention. \textit{Travelers' health.} https://wwwnc.cdc.gov/travel/
    \item NHS Fit for Travel. \textit{Travel vaccination advice.} https://www.fitfortravel.nhs.uk/
\end{itemize}
